\section{Related Work}
\subsection{Clinical datasets, representation and shortcuts}
Dermatology collections have enabled transfer learning but have also exposed variability in annotation and acquisition source. Fitzpatrick17k uses clinical photographs annotated for skin type, highlighting the importance of representation in evaluating learning systems \citep{groh2021fitzpatrick}. SCIN broadens the setting through user-contributed images and condition information \citep{ward2024scin,scin_repository}. Neither removes the need to check anatomical site, original diagnosis, photographic modality and permitted use when assembling a new task. Different crops of a single photograph, watermarks or site-specific presentation can also create shortcuts. If the same visual subject reaches both training and testing, a classifier may appear to generalize while recognizing an already-seen image. We therefore consider label and leakage auditing part of the research design, not merely data cleaning.

\subsection{Diffusion-generated dermatology images}
Latent diffusion offers a flexible approach to text-conditioned synthesis \citep{rombach2022latent}. Dermatology research has begun exploring whether generated lesions help train classifiers; Derm-T2IM studies synthetic \emph{dermoscopic} lesions with vision transformers and CNNs \citep{farooq2024dermt2im}. A dermoscopic crop differs from full-face photography in its anatomical context, image scale and acquisition artifacts. Results from one modality are not evidence that generated faces improve another. Moreover, a model may synthesize a plausible skin texture while omitting the requested lesion, distorting its border or introducing a consistent synthetic artifact. In this study, lesion fidelity, image novelty and classification utility are deliberately measured separately. An aesthetic score or embedding similarity alone is not treated as a clinical truth label.

\subsection{General and domain-specialized encoders}
ResNet-50 supplies residual convolutional learning \citep{he2016resnet}; EfficientNet-B3 represents compound scaling \citep{tan2019efficientnet}; and ConvNeXt Tiny modernizes convolutional blocks while remaining a CNN \citep{liu2022convnext}. DINOv2 ViT-S/14 provides features learned by large-scale self-supervision \citep{oquab2024dinov2}. These families offer distinct feature representations for testing whether an augmentation effect transfers across classifiers rather than appearing on one backbone only.

PanDerm is especially relevant as a dermatology-pretrained foundation model trained across multiple imaging modalities \citep{yan2025panderm}. Its published results motivate examination of specialty representations but do not establish superiority on our six-category facial task. Our PanDerm Base experiment is a frozen-feature linear probe; it must not be described as full fine-tuning. The four general-purpose backbones instead support head training, final-stage unfreezing and optional full fine-tuning.

\subsection{Unresolved methodological question}
The gap addressed here is not the absence of synthetic dermatology work. It is the difficulty of attributing an improvement to a specific intervention if prompt detail, generator weights, synthetic-image selection, classifier optimization and test composition all change simultaneously. DermSynReal separates the augmentation, prompt and selection effects into explicit comparisons. When an intended match is not feasible, the result is described as a confounded exploratory comparison rather than evidence that morphology wording caused the difference.
